\subsection{Supervision within Data}
\textbf{Language supervision} \citet{joulin2016learning, gomez2017self, zhang2020contrastive, sariyildiz2020learning, desai2021virtex} demonstrate the effectiveness of learning transferable visual features from language supervision. 
% However, these works are limited to a small dataset such as Flickr~\citep{young2014image} or COCO~\citep{lin2014microsoft}. 
Pioneering work CLIP~\citep{radford2021learning} and ALIGN~\citep{jia2021scaling} achieve prestigious performance via learning from 400M/1B image-text pairs. We are following these two works to improve their data efficiency.

Relevant concurrent work including: SLIP~\cite{mu2021slip} introduces self-supervision to CLIP. 
FILIP~\cite{yao2021filip} leverages the finer-grained alignment between image patches and textual words.
OTTER~\cite{wu2021data,cheng2021data} uses online entropic optimal transport to find a soft image-text match as labels to mitigate the noise within the dataset. 

% They only use the simple image-text pair contrastive supervision while ignoring other widespread supervisions among the pairs, thereby are inefficient.

\textbf{Visual self-supervision} Our work is also highly related to self-supervised learning (SSL)~\citep{lecun2021self}. 
Contrastive learning, as a pretext task of SSL, has achieved remarkable success in visual representation learning~\citep{he2020momentum, chen2020simple, caron2020unsupervised, grill2020bootstrap, chen2020simple}. Researchers also extend contrastive learning into multi-modal settings~\citep{yuan2021multimodal}. However, it is only limited to a small COCO dataset~\citep{yuan2021multimodal}.
% improve SSL, Multi-modal contrastive learning~\citep{yuan2021multimodal}
% To further improve SSL, researchers either embrace intra-modal supervision~\citep{yuan2021multimodal}. However, they are limited to a small COCO dataset~\citep{yuan2021multimodal}.

\textbf{Nearest-neighbor supervision} 
Recently, researchers have exploited nearest-neighbor supervision to learn visual features~\citep{dwibedi2021little, van2021revisiting}.
% the image nearest-neighbor supervision is proposed as a new self-supervised learning paradigm for representative image feature extraction~\citep{dwibedi2021little, van2021revisiting}.
They find that using nearest-neighbor as positive samples in the contrastive loss improves the performances on multiple downstream tasks. However, they mainly focus on the single visual modality pretraining on relatively small datasets, such as ImageNet. We propose novel nearest-neighbor supervision for multi-modal learning to harness information from other similar pairs.

% To our best view, our \workname is the first work to dig into the supervisory signals among the million-scale image-text dataset. 
% Our work opens up a new direction to fully exploit the existing multi-modal data instead of scaling up data naively. 
% Moreover, o
% these approaches learn representations by contrasting positive pairs against negative pairs.
% There are also literature that focus on improving the quli via embracing of multimodal data~\citep{yuan2021multimodal} or using nearest neighborhood.
% Benefiting from the million scale image-text pairs,

% Deep neural networks, such as CNNs~\citep{krizhevsky2012imagenet, he2016deep} or transformers~\citep{vaswani2017attention, dosovitskiy2020image}, can automatically learn discriminative feature representations from data according to the supervisory signals. The supervisions can either be external human labeling, a.k.a supervised learning~\citep{lecun2015deep}, or internal structural information within the data, a.k.a self-supervised learning~\citep{lecun2021self} (SSL). 


